\documentclass[11pt]{article}
\usepackage[margin=0.55in]{geometry}
\usepackage{amsmath,amssymb}
\usepackage{xcolor}
\usepackage{tcolorbox}
\usepackage{lmodern}
\usepackage[T1]{fontenc}
\pagestyle{empty}

\definecolor{hkustblue}{RGB}{0,56,116}
\definecolor{lightblue}{RGB}{235,244,255}
\definecolor{lightgray}{RGB}{248,248,248}

\newtcolorbox{keybox}{
  colback=lightblue,
  colframe=hkustblue,
  boxrule=0.8pt,
  arc=2pt,
  left=6pt,
  right=6pt,
  top=6pt,
  bottom=6pt
}

\begin{document}

\begin{center}
{\Large \bfseries HW2 Q3: Capacity Allocation after Poisson Splitting}\\[0.2cm]
{\normalsize IEDA 1250 Queueing Theory}
\end{center}

\vspace{0.1cm}

\textbf{Can we treat the system as two independent M/M/1 queues?}

Yes. By Poisson splitting, the two arrival streams are Poisson with rates
\[
\lambda_1=p\lambda,\qquad \lambda_2=(1-p)\lambda.
\]
So we can model the two branches as two M/M/1 queues, but the overall mean response
time must be weighted by the routing probabilities:
\[
E[T]
=
p\frac{1}{\mu_1-p\lambda}
+
(1-p)\frac{1}{\mu_2-(1-p)\lambda},
\qquad
\mu_1+\mu_2=\mu.
\]

\begin{keybox}
\textbf{Stability conditions}
\[
\mu_1>p\lambda,\qquad
\mu_2>(1-p)\lambda.
\]
These are the minimum capacities needed just to keep both queues stable.
\end{keybox}

\textbf{Easy case: \(p=\frac{1}{2}\).}

By symmetry, split capacity equally:
\[
\boxed{\mu_1^*=\mu_2^*=\frac{\mu}{2}.}
\]
Each queue has arrival rate \(\lambda/2\), so
\[
E[T_1]=E[T_2]
=\frac{1}{\mu/2-\lambda/2}
=\frac{2}{\mu-\lambda}.
\]
Thus
\[
\boxed{E[T]=\frac{2}{\mu-\lambda}.}
\]

\textbf{General case: \(p>\frac{1}{2}\).}

First allocate the minimum stable capacities:
\[
p\lambda \quad \text{to server 1},\qquad
(1-p)\lambda \quad \text{to server 2}.
\]
The remaining extra capacity is
\[
\mu-\lambda.
\]
Define the extra capacities
\[
x_1=\mu_1-p\lambda,\qquad
x_2=\mu_2-(1-p)\lambda.
\]
Then
\[
x_1+x_2=\mu-\lambda,
\qquad
E[T]=\frac{p}{x_1}+\frac{1-p}{x_2}.
\]
Minimizing this gives
\[
\frac{x_1}{x_2}
=
\frac{\sqrt{p}}{\sqrt{1-p}}.
\]

\begin{keybox}
\textbf{Optimal fraction of extra capacity going to server 1}
\[
\boxed{
\frac{x_1}{x_1+x_2}
=
\frac{\sqrt{p}}{\sqrt{p}+\sqrt{1-p}}
}
\]
\end{keybox}

Therefore, the optimal capacity allocation is
\[
\boxed{
\mu_1^*
=
p\lambda
+
(\mu-\lambda)
\frac{\sqrt{p}}{\sqrt{p}+\sqrt{1-p}}
}
\]
\[
\boxed{
\mu_2^*
=
(1-p)\lambda
+
(\mu-\lambda)
\frac{\sqrt{1-p}}{\sqrt{p}+\sqrt{1-p}}
}
\]

\vspace{0.1cm}
\textbf{Intuition:} each server first receives enough capacity to handle its own arrival
rate. The remaining capacity should be split in proportion to
\(\sqrt{p}\) and \(\sqrt{1-p}\), not directly in proportion to \(p\) and \(1-p\).

\end{document}
